% Nouveau EGA V (cote 25): archive cote 25, site page 142 (PDF page 143).
% typescript. Draft compared: none.
% AI-transcribed by Claude Opus 5.5 (claude-opus-5-5, adaptive thinking, effort high), Message Batches API, batch ega25b (msgbatch_01A7n5YuXubUPNK1zsepi8RG), 2026-11-13; not yet checked by eye.
\scanpage{25/142}
\mspage{\hand{3}}

préschéma x sur k, \var{\xout{enfin la dernière assertion résulte de la première}}{}{typewriter strike-out}{ega25_25_142-v0001}
\var{\xout{précédente et de l'isomorphisme $\underline{m}_x$ analogue $gr^2(P^2_{X/k})(x)$ Sy}}{}{typewriter strike-out}{ega25_25_142-v0002}

Notons que $gr^1(\underline{P}^1_{X/Y}) \hand{\simeq} \underline{\hand{\Omega}}^1_{X/Y}$ , par suite la restriction
$d^1\hand{\phi}\,|V(\hand{\phi})$ peut être interprété comme une section de $\underline{\hand{\Omega}}^1_{X/Y} \otimes \underline{O}_{V(\hand{\phi})}$ ,
qui n'est autre d'ailleurs que la restriction de $d_{X/Y}\,\hand{\phi}$ à $V(\hand{\phi})$ . Nous
pouvons alors considérer le \uline{préschéma des zéros de cette section} , que
nous noterons $V(\hand{\phi})^{sing}$ , et dont l'ensemble sous-jacent n'est autre
que l'ensemble des zéros de $\hand{\phi}$ singuliers rel à Y en vertu de prop.5 .
( NB si \var{\xout{\unsure{$\hand{\phi}$}}}{}{typewriter strike-out}{ega25_25_142-v0003} $\hand{\mu}$ est une section d'un Module localement libre \var{\inserted[typed]{de type fini}}{de type fini}{insertion}{ega25_25_142-v0004} $\underline{E}$ sur un
\var{\margin{$\|$}}{\placedmargin{$\|$}}{marginal note placed}{ega25_25_142-v0005}\var{\inserted[typed]{préschéma}}{préschéma}{insertion}{ega25_25_142-v0006} \var{\xout{espace annelé}}{}{typewriter strike-out}{ega25_25_142-v0007} X, on définit de façon évidente le sous-préschéma d es
zéros de $\hand{\mu}$ , par exemple comme défini par l'Idéal image de $\overset{\hand{\scriptscriptstyle\vee}}{\underline{E}} \hand{\longrightarrow} \underline{O}_X$
transposé de $\hand{\mu}$ ; lorsque $\underline{E} = \underline{O}_X{}^n$ et $\hand{\mu} = \var{\xout{\unsure{(\hand{\mu}_i)}}}{}{typewriter strike-out}{ega25_25_142-v0008}\,(\hand{\mu}_1,..,\hand{\mu}_n)$, alors cet
idéal n'est autre que $\hand{\sum} \hand{\mu}_i\,\underline{O}_X$ ,qui définit $V(\hand{\mu}_1,..,\hand{\mu}_n)$ ). Prenant
maintenant la restriction $d^2\hand{\phi}\,|V(\hand{\phi})^{sing}$ , et notant que
$gr^2(P^2_{X/Y}) \hand{\simeq} Sym\var{\xout{m}}{}{typewriter strike-out}{ega25_25_142-v0009}^2(\underline{\hand{\Omega}}^1_{X/Y})$, on trouve une section canonique $M(\hand{\phi})$ de
$Sym^2(\underline{\hand{\Omega}}^{\,1}_{X/\unsure{Y|X}}) \otimes \underline{O}_{V^{sing}}$ . On vérifie immédiatement que (prenant des
\var{\margin{$\|$}}{\placedmargin{$\|$}}{marginal note placed}{ega25_25_142-v0010}point de X à val\var{\xout{eu}}{}{typewriter strike-out}{ega25_25_142-v0011}rs dans des corps ...) cette section est précisément
celle qui détermine les formes quadratiques envisagées dans déf 1
(dans le cas \var{\xout{\illegible{1}}}{}{typewriter strike-out}{ega25_25_142-v0012} d'un $X_k$ déduit de X/S par Spec(k)$\hand{\rightarrow}$ S). On
en déduit une description d e l'ensemble $V(\hand{\phi})^{sup\ sing}$ en termes de
cette section, de la façon suivante: interprétant $M(\hand{\phi})$ comme définissan
% NOTE: the end of this line is cut by the right edge of the scan (presumably "définissant" / "définissant un")
homomorphisme
\[
M(\hand{\phi})\hand{'} : \ \hand{\mathcal{T}}_{X/Y} \otimes_{\unsure{\underline{O}_X}} \unsure{\underline{O}_{V^{sing}}} \hand{\longrightarrow} \underline{\hand{\Omega}}^1_{X/Y} \otimes_{\underline{O}_X} \underline{O}_{V^{sing}} \quad ,
\]
il faut prendre l'ensemble des points en lesquels cet homomorphisme
n'est pas un isomorphisme. Cela montre en particulier que $V(\hand{\phi})^{sup\ sing}$
est un ensemble fermé. On peut \var{\deleted{le}}{}{deletion}{ega25_25_142-v0013} préciser \var{\inserted{ce dernier}}{ce dernier}{insertion}{ega25_25_142-v0014}\var{\deleted{, en}}{}{deletion}{ega25_25_142-v0015} introduisant
\[
D(\hand{\phi}) = \text{dét}\ M(\hand{\phi}) \ \hand{\in} \ \hand{\Gamma}\,(\underline{\hand{\Omega}}^d_{X/Y}{}^{\otimes 2} \otimes \underline{O}_{V^{sing}}\ ) \quad ,
\]

\mspage{\hand{142}}
